%% ma: L10-B01-0023
%% lop: 10
%% chuong: 1
%% bai: 1
%% dang: 10.1.6
%% loai: TN4
%% muc: TH
%% dap_an: "A"
%% nguon: "Ảnh giáo viên gửi 10/10/2026 (ThS. Nguyễn Hoàng Việt, Chương 1 Mệnh đề), B-Đề 2, Phần 1 Câu 3"
%% kien_thuc: "Bài 1 (mệnh đề kéo theo, mệnh đề tương đương)"
%% trang_thai: cho-duyet
%% ly_do: "Dạng toán mới đề xuất 10.1.6: mệnh đề kéo theo, tương đương: xét tính đúng sai"
\begin{ex}
Trong các mệnh đề nào sau đây mệnh đề nào \textbf{sai}?
\choice{\True Hai tam giác bằng nhau khi và chỉ khi chúng đồng dạng và có một góc bằng nhau}{Một tứ giác là hình chữ nhật khi và chỉ khi chúng có $3$ góc vuông}{Một tam giác là vuông khi và chỉ khi nó có một góc bằng tổng hai góc còn lại}{Một tam giác là đều khi và chỉ khi chúng có hai đường trung tuyến bằng nhau và có một góc bằng $60^\circ$}
\loigiai{A sai: hai tam giác đồng dạng theo tỉ số $2$ có góc bằng nhau nhưng không bằng nhau. B đúng. C đúng vì tổng ba góc bằng $180^\circ$ nên góc đó bằng $90^\circ$. D đúng: hai trung tuyến bằng nhau thì tam giác cân; cân và có một góc $60^\circ$ thì đều.}
\end{ex}
